\begin{table*}[t]
  \centering\small
  \begin{tabular}{@{}llrrrrrr@{}}
    \toprule
    Rows & Frailty & Spells & Deaths & Slope & 95\% CI & $\theta$ & Absorbed \\
    \midrule
    all spells & no frailty (baseline) & 247,694 & 28,255 & $-0.0506$ & $[-0.0524, -0.0489]$ & -- & -- \\
     & shared within repository & 247,694 & 28,255 & $-0.0481$ & $[-0.0499, -0.0464]$ & 1.44 & 4.9\% \\
     & shared within file & 247,694 & 28,255 & $-0.0485$ & $[-0.0502, -0.0467]$ & 1.49 & 4.3\% \\
    recurring text & no frailty (baseline) & 20,807 & 2,669 & $-0.0513$ & $[-0.0571, -0.0455]$ & -- & -- \\
     & shared within repository & 20,807 & 2,669 & $-0.0505$ & $[-0.0566, -0.0445]$ & 2.85 & 1.5\% \\
     & shared by instruction text & 20,807 & 2,669 & $-0.0355$ & $[-0.0414, -0.0296]$ & 3.41 & 30.8\% \\
    \bottomrule
  \end{tabular}
  \caption{
    Composition shrinks the age slope at every clustering level and never reverses it.
    Each row is one gamma-frailty specification, fit by piecewise-exponential EM
    (\S\ref{sec:appendix:frailty}). The columns give the log-hazard slope per commit with its 95\%
    CI, the fitted frailty variance $\theta$, and the share of that block's no-frailty slope which
    the clustering absorbs. \emph{Recurring text} restricts to instructions whose text appears in
    $\ge 2$ repositories, which are the rows where a frailty shared by instruction content can be
    identified; absorbed shares compare within a block, against the no-frailty fit on those same
    rows. We censor deaths administratively at 50 commits. A slope of $0$ would put deletion risk
    flat in instruction age, and a positive slope is what instruction staleness predicts. Content
    absorbs the most. Repository clustering on those same rows absorbs almost nothing, which puts
    the absorbed variance in the instruction text rather than in its author.
  }
  \label{tab:frailty}
\end{table*}
